\chapter{Swaps de variance et produits dérivés de volatilité}\label{ch:dv:variance-swaps-and-volatility-derivatives}

Un \omterm{def:dv:variance-swaps-and-volatility-derivatives:vs}{swap de variance} n'a pas de strike exprimé en prix. Son seul strike est une volatilité, élevée au carré, et il verse la différence
entre ce nombre et la variance que l'indice réalise effectivement. Il ressemble à un pari qui exige un modèle de
volatilité. Il n'en exige aucun. Sa couverture est un strip statique de toutes les options hors de la monnaie sur l'indice, chacune
pondérée par l'inverse de son strike au carré, plus une position sur l'indice rééquilibrée chaque jour. L'indice de
volatilité est le prix de ce strip, coté en volatilité. Ce chapitre construit la réplication, la formule de l'indice et ses
erreurs, ainsi que la convexité qui distingue un swap sur volatilité d'un swap sur variance.
Il traite ensuite des contrats à terme et des options sur l'indice de volatilité, ainsi que des sauts et des strikes tronqués qui
rompent la réplication.

\section{Swaps de variance et leur réplication}

\begin{definition}[Swap de variance]\label{def:dv:variance-swaps-and-volatility-derivatives:vs}
Un \emph{swap de variance}\index{swap de variance} est un contrat à terme sur la variance réalisée. À l'échéance $T$, le
côté acheteur reçoit $N_{\mathrm{var}}\bigl(\sigma_R^2-K^2\bigr)$, où
$\sigma_R^2=\frac{252}{n}\sum_{i=1}^n\ln^2(S_{t_i}/S_{t_{i-1}})$ est la variance réalisée annualisée des cours de clôture quotidiens, sans
soustraction de moyenne, et $K$ est le strike, coté en volatilité.
\end{definition}

\begin{definition}[Notionnel de variance]\label{def:dv:variance-swaps-and-volatility-derivatives:varnot}
Le \emph{notionnel de variance}\index{notionnel de variance} $N_{\mathrm{var}}$ est le paiement par unité de variance
(par «~point de variance~» lorsque la variance est exprimée en pourcentage au carré).
\end{definition}

\begin{definition}[Notionnel vega]\label{def:dv:variance-swaps-and-volatility-derivatives:veganot}
Le \emph{notionnel vega}\index{notionnel vega} d'un \omterm{def:dv:variance-swaps-and-volatility-derivatives:vs}{swap de variance} est $N_{\mathrm{vega}}=2KN_{\mathrm{var}}$~: le
paiement approximatif pour un point de volatilité réalisée au-dessus du strike, au voisinage du strike.
Les transactions sont cotées en notionnel \omterm{def:dv:greeks-and-the-hedging-pnl:greeks}{vega} et réglées en \omterm{def:dv:variance-swaps-and-volatility-derivatives:varnot}{notionnel de variance}.
\end{definition}

Le paiement est convexe en volatilité réalisée. Avec un \omterm{def:dv:variance-swaps-and-volatility-derivatives:veganot}{notionnel vega} de 100\,000 par point et un strike de
22.1~\%, le \omterm{def:dv:variance-swaps-and-volatility-derivatives:varnot}{notionnel de variance} est $100\,000/(2\times22.1)=2\,262$ par point de variance. Une volatilité réalisée
de 30~\% verse $2\,262\times(30^2-22.1^2)=931\,000$ au lieu des 790\,000 qu'un paiement linéaire donnerait, et
15~\% coûte 596\,000 au lieu de 710\,000. L'acheteur gagne plus que proportionnellement lorsque la volatilité fait un bond.
Cette convexité explique pourquoi un \omterm{def:dv:variance-swaps-and-volatility-derivatives:vs}{swap de variance} est réplicable et un \omterm{def:dv:variance-swaps-and-volatility-derivatives:volswap}{swap de volatilité} ne l'est pas.

\begin{proposition}[Réplication]\label{prop:dv:variance-swaps-and-volatility-derivatives:rep}
Si l'indice évolue de manière continue, quel que soit le processus de volatilité, alors
\[
\int_0^T\sigma_t^2\,dt=2\int_0^T\frac{dS_t}{S_t}-2\ln\frac{S_T}{S_0},
\]
et, avec $F$ le forward et les taux zéro-coupon,
\[
-2\ln\frac{S_T}{F}=-2\,\frac{S_T-F}F+\int_0^F\frac2{K^2}(K-S_T)^+\,dK+\int_F^\infty\frac2{K^2}(S_T-K)^+\,dK.
\]
La variance réalisée est donc le paiement d'un strip statique d'options hors de la monnaie, pondérées
par $2/K^2$, d'un forward, et d'une position dynamique de $2/S_t$ sur l'indice. Le strike équitable est
\[
K^2T=2\int_0^F\frac{P(K)}{K^2}\,dK+2\int_F^\infty\frac{C(K)}{K^2}\,dK.
\]
\end{proposition}
\begin{proof}
La formule d'Itô (livre~4, chapitre~3) pour $\ln S$ donne
$d\ln S_t=dS_t/S_t-\frac12\sigma_t^2\,dt$~; on intègre et on réarrange. Pour la seconde identité, toute fonction de paiement deux fois
dérivable vérifie
$f(S)=f(F)+f'(F)(S-F)+\int_0^Ff^{\prime\prime}(K)(K-S)^+dK+\int_F^\infty f^{\prime\prime}(K)(S-K)^+dK$ (formule de Taylor
à reste intégral), et $f(S)=-2\ln(S/F)$ vérifie $f^{\prime\prime}(K)=2/K^2$. La composante dynamique $2\int dS/S$ a un prix nul, et il en est de même
du forward, ce qui ne laisse que le strip.
\end{proof}

La démonstration n'a jamais utilisé de modèle pour $\sigma_t$. C'est là l'attrait du produit~: le strike se lit dans les prix
des seules options. Le dealer qui vend le swap achète le strip une fois, négocie l'indice chaque jour et est
couvert, quel que soit le comportement de la volatilité, tant que l'indice ne fait pas de saut. Les poids du strip indiquent aussi
de quoi le strike est fait (\cref{fig:dv:variance-swaps-and-volatility-derivatives:strip}). Sur la surface du
chapitre~9, le strike à un an est de 22.1~\%, contre une volatilité à la monnaie de 19.2~\%. Les puts portent 72~\% de la
variance, les strikes inférieurs à 80 en portent 24~\%, et ceux supérieurs à 120 en portent 1.5~\%. Un \omterm{def:dv:variance-swaps-and-volatility-derivatives:vs}{swap de variance} sur un indice présentant un skew est
pour l'essentiel une position sur sa partie baissière.

\begin{omfigure}
\begin{tikzpicture}
\begin{groupplot}[group style={group size=2 by 1,horizontal sep=1.6cm},
  width=6.6cm,height=5cm,
  tick label style={font=\small},label style={font=\small},grid=major,grid style={gray!25},table/col sep=comma]
\nextgroupplot[xlabel={strike},ylabel={part de la variance (\%)},xmin=25,xmax=165,ymin=0,ymax=17,
  title={\small strip à un an, strikes tous les 5}]
\addplot[omThm,ybar,bar width=3.5pt,fill=omThm!40] table[x=k,y=share]{figdata/derivatives/14-variance-swaps-and-volatility-derivatives/contrib.csv};
\draw[gray,dashed] (axis cs:100,0) -- (axis cs:100,17);
\nextgroupplot[xlabel={indice à l'échéance},ylabel={paiement},xmin=40,xmax=170,ymin=-1.2,ymax=1.9,
  title={\small $-2\ln(S_T/F)$ et un strip},
  legend columns=1,legend cell align=left,legend style={font=\footnotesize,at={(0.5,-0.32)},anchor=north,draw=none,fill=none}]
\addplot[omDef,very thick] table[x=s,y=log]{figdata/derivatives/14-variance-swaps-and-volatility-derivatives/replication.csv};
\addplot[omThm,thick,dashed] table[x=s,y=strip]{figdata/derivatives/14-variance-swaps-and-volatility-derivatives/replication.csv};
\legend{paiement logarithmique,{strip, strikes de 60 à 140 tous les 10}}
\end{groupplot}
\end{tikzpicture}

\omcaption{À gauche~: la part de chaque strike dans le strike du \omterm{def:dv:variance-swaps-and-volatility-derivatives:vs}{swap de variance} à un an sur la surface du chapitre~9~: les puts
situés sous le forward (en pointillés) portent 72~\%. À droite~: le paiement logarithmique et sa réplication par neuf strikes. Le strip est
linéaire par morceaux, exact aux strikes, à l'erreur d'espacement près, et linéaire au-delà des derniers strikes, où il
cesse de répliquer. Données~: le code du chapitre.}\label{fig:dv:variance-swaps-and-volatility-derivatives:strip}
\end{omfigure}

Un \omterm{def:dv:variance-swaps-and-volatility-derivatives:vs}{swap de variance} fait l'objet d'une valorisation au prix de marché par découpage dans le temps. Après un instant $t$ avec une variance réalisée
$\sigma_{R,t}^2$, la variance finale attendue est la moyenne pondérée par le temps
$\bigl(t\,\sigma_{R,t}^2+(T-t)K_{t,T}^2\bigr)/T$ de la partie réalisée et du strike équitable de la
partie restante. La valeur vaut $N_{\mathrm{var}}$ fois l'écart actualisé par rapport à $K^2$. Un swap à un an fixé à un strike de
22.1~\%, ayant réalisé 25~\% sur six mois, avec une volatilité de \omterm{def:dv:variance-swaps-and-volatility-derivatives:vs}{swap de variance} à six mois de 20.1~\%, attend
$0.5\times25^2+0.5\times20.1^2=514.5$ points de variance contre un strike de 488.4. Sur un \omterm{def:dv:variance-swaps-and-volatility-derivatives:varnot}{notionnel de variance} de
2\,262, il vaut 59\,000.

\section{Le contrat log et la formule de l'indice de volatilité}

\begin{definition}[Contrat log]\label{def:dv:variance-swaps-and-volatility-derivatives:log}
Le \emph{contrat log}\index{contrat log} verse $\ln(S_T/F)$ à $T$. Son prix, $-\frac12\sigma_{\mathrm{VS}}^2T$
avec des taux nuls, est le strike du \omterm{def:dv:variance-swaps-and-volatility-derivatives:vs}{swap de variance} exprimé dans une autre unité. Il n'est pas coté et se négocie comme son
strip d'options.
\end{definition}

L'indice de volatilité transforme le strip en un chiffre publié. Sur une chaîne d'options cotées, l'intégrale devient une
somme sur les strikes, et la formule de l'indice est
\[
\sigma^2=\frac2T\sum_i\frac{\Delta K_i}{K_i^2}\,e^{RT}Q(K_i)-\frac1T\Bigl(\frac F{K_0}-1\Bigr)^2,
\]
avec $Q(K_i)$ le cours médian de l'option hors de la monnaie en $K_i$, $K_0$ le strike égal ou immédiatement inférieur au
forward, et $\Delta K_i$ la moitié de la distance entre strikes voisins. Le dernier terme corrige le fait d'utiliser
le put et le call en $K_0$ plutôt qu'en $F$. Entre $K_0$ et $F$, le strip détient un call dans la monnaie sous la forme
d'un put hors de la monnaie, et le terme élimine l'écart à l'ordre deux. Deux échéances encadrant trente
jours sont calculées et interpolées à une durée constante de trente jours.

\begin{dated}{2026-09}{La méthodologie de l'indice de volatilité}\label{dat:dv:variance-swaps-and-volatility-derivatives:vix}
Le document de méthodologie de Cboe pour cet indice (version 6.0, révisé pour la dernière fois le 26 février 2026) précise la formule ci-dessus.
Il utilise les options sur le S\&P~500 au cours médian, ne retient que les options à cours acheteur non nul, encadre trente jours avec deux
échéances, standard ou hebdomadaires, et désigne comme fondement de la formule la note de recherche de 1999 de Demeterfi, Derman,
Kamal et Zou. Un algorithme de filtrage protège contre les cotations dont l'écart s'élargit brusquement. Les contrats à terme et les options sur
l'indice règlent sur un cours spécial d'ouverture (SOQ) qui suit des règles légèrement différentes de celles de l'indice au comptant.
\end{dated}

Quelle est la qualité de la somme~? Elle comporte deux erreurs~: l'espacement entre les strikes, et la plage couverte par les strikes, qui, dans
l'indice, est tronquée là où les cours acheteurs disparaissent. Sur le smile à trente jours du chapitre~9, dont la volatilité exacte du strip vaut 16.30~\%,
le chiffre de type indice présente un écart comme le montre le \cref{tab:dv:variance-swaps-and-volatility-derivatives:errors}. Un espacement large
la biaise à la hausse, car une somme grossière sur une fonction convexe la surestime. Une plage courte la biaise à la baisse, puisque les
ailes manquantes auraient ajouté de la variance. Un strip qui s'arrête à 95 manque 1.6 à 1.8 points. Les deux erreurs peuvent
se compenser par hasard, comme le montre le strip partant de 90 avec un pas de 2.5.

\begin{table}[ht]
\centering
\begin{tabular}{rrrr}
\toprule
strike le plus bas (le plus haut $=200-$ strike le plus bas) & $\Delta K=0.5$ & $\Delta K=1$ & $\Delta K=2.5$\\
\midrule
70 & $+0.01$ & $+0.06$ & $+0.38$\\
80 & $-0.02$ & $+0.03$ & $+0.36$\\
85 & $-0.13$ & $-0.07$ & $+0.28$\\
90 & $-0.48$ & $-0.40$ & $+0.01$\\
95 & $-1.78$ & $-1.61$ & $-0.92$\\
\bottomrule
\end{tabular}
\caption{Erreur de la volatilité à trente jours de type indice par rapport au strip exact (16.30~\%), en points de volatilité,
selon la plage de strikes et l'espacement, sur la surface du chapitre~9 (spot 100). Données~: le tutoriel.}\label{tab:dv:variance-swaps-and-volatility-derivatives:errors}
\end{table}

\section{Swaps de volatilité et convexité}

\begin{definition}[Swap de volatilité]\label{def:dv:variance-swaps-and-volatility-derivatives:volswap}
Un \emph{swap de volatilité}\index{swap de volatilité} verse $N\bigl(\sigma_R-K_{\mathrm{vol}}\bigr)$~: la volatilité
réalisée, et non la variance, contre un strike.
\end{definition}

Le paiement est linéaire en volatilité, ce que beaucoup de clients veulent négocier, mais la volatilité est la racine
carrée de la variance, et aucun strip statique ne réplique la racine carrée d'une intégrale sur le chemin. Son strike équitable est
$\E[\sigma_R]$, et l'inégalité de Jensen la place sous le strike du \omterm{def:dv:variance-swaps-and-volatility-derivatives:vs}{swap de variance} $\sqrt{\E[\sigma_R^2]}$. Un
développement à l'ordre deux de la racine carrée autour de la variance moyenne $V$ donne
\[
K_{\mathrm{vol}}\approx\sqrt V-\frac{\Var(\sigma_R^2)}{8V^{3/2}},
\]
l'ajustement de convexité est donc proportionnel à la variance de la variance réalisée~: il faut un modèle pour la
\omterm{def:dv:stochastic-volatility:sv}{volatilité de la volatilité}. Il s'agit de l'ajustement de convexité du livre~2, chapitre~8, appliqué à une
racine carrée.

Dans le \omterm{def:dv:stochastic-volatility:heston}{modèle de Heston}, il est exact. La transformée de Laplace de la variance intégrée,
$\E[e^{-s\int_0^Tv_t\,dt}]$, admet la forme fermée d'un modèle affine. L'identité
$\sqrt x=\frac1{2\sqrt\pi}\int_0^\infty(1-e^{-sx})s^{-3/2}ds$ la transforme en $\E[\sqrt{\cdot}]$ par une
seule intégrale. La construction fait exactement cela, et ses tests la confrontent à une simulation.

\begin{example}[Volatilité ou variance]\label{ex:dv:variance-swaps-and-volatility-derivatives:convexity}
Dans le \omterm{def:dv:stochastic-volatility:heston}{modèle de Heston} calibré sur la surface du chapitre~9 (chapitre~10), le strike du \omterm{def:dv:variance-swaps-and-volatility-derivatives:vs}{swap de variance} à un an est
21.4~\% et le strike du \omterm{def:dv:variance-swaps-and-volatility-derivatives:volswap}{swap de volatilité} 19.8~\%~: un ajustement de convexité de 1.62 points de volatilité. L'ajustement
vaut 0.72 point à un mois, 1.34 à trois mois et 1.61 à six mois, puis tombe à 1.35 à deux ans et 1.10
à trois ans (la \cref{fig:dv:variance-swaps-and-volatility-derivatives:convexity}). Aux échéances courtes, il reste peu
de temps pour que la variance varie. Aux échéances longues, le retour à la moyenne en lisse la variation. Le strip de marché
à un an donne 22.1~\%, 0.7 point au-dessus de celui de Heston, parce que le modèle calibré ne reproduit pas
les ailes de la surface, qui pèsent dans le strip.
\end{example}

\begin{omfigure}
\begin{tikzpicture}
\begin{axis}[width=10.5cm,height=5cm,
  xlabel={échéance (années)},ylabel={strike (volatilité, \%)},
  tick label style={font=\small},label style={font=\small},
  xmin=0,xmax=3.1,ymin=15,ymax=25,grid=major,grid style={gray!25},
  legend columns=3,legend cell align=left,legend style={font=\footnotesize,at={(0.5,-0.3)},anchor=north,draw=none,fill=none,/tikz/every even column/.append style={column sep=8pt}},
  table/col sep=comma]
\addplot[omDef,very thick,mark=*,mark size=1.6pt] table[x=t,y=var]{figdata/derivatives/14-variance-swaps-and-volatility-derivatives/heston_terms.csv};
\addplot[omThm,very thick,mark=square*,mark size=1.6pt] table[x=t,y=vol]{figdata/derivatives/14-variance-swaps-and-volatility-derivatives/heston_terms.csv};
\addplot[gray,thick,dashed,mark=o,mark size=1.6pt] table[x=t,y=market]{figdata/derivatives/14-variance-swaps-and-volatility-derivatives/heston_terms.csv};
\legend{swap de variance (Heston),swap de volatilité (Heston),swap de variance (strip de marché)}
\end{axis}
\end{tikzpicture}

\omcaption{Strikes des \omterm{def:dv:variance-swaps-and-volatility-derivatives:vs}{swaps de variance} et des \omterm{def:dv:variance-swaps-and-volatility-derivatives:volswap}{swaps de volatilité} par échéance dans le \omterm{def:dv:stochastic-volatility:heston}{modèle de Heston} calibré sur la
surface du chapitre~9, et strike du \omterm{def:dv:variance-swaps-and-volatility-derivatives:vs}{swap de variance} du strip de marché. L'écart entre les deux premières séries est
l'ajustement de convexité, maximal vers un an. Données~: le tutoriel.}\label{fig:dv:variance-swaps-and-volatility-derivatives:convexity}
\end{omfigure}

La variance peut être pondérée et restreinte, et deux variantes sont négociées. Chacune est répliquée par le même
raisonnement, avec une fonction différente du prix.

\begin{definition}[Gamma swap]\label{def:dv:variance-swaps-and-volatility-derivatives:gamma}
Un \emph{gamma swap}\index{gamma swap} verse la variance réalisée pondérée par le niveau de l'indice rapporté à sa
valeur initiale, $\frac1T\int_0^T(S_t/S_0)\,\sigma_t^2\,dt$, contre un strike. Il est répliqué par le paiement
$\frac2{S_0}(S\ln(S/S_0)-S+S_0)$, dont la dérivée seconde vaut $2/(S_0K)$~: son strip pondère les puts moins que
celui du \omterm{def:dv:variance-swaps-and-volatility-derivatives:vs}{swap de variance}.
\end{definition}

\begin{definition}[Swap de variance corridor]\label{def:dv:variance-swaps-and-volatility-derivatives:corridor}
Un \emph{swap de variance corridor}\index{swap de variance corridor} accumule la variance réalisée seulement les jours où
l'indice se situe dans un intervalle $[L,U]$. Son strip ne conserve les pondérations $2/K^2$ que sur les strikes de $[L,U]$.
\end{definition}

La variance du \omterm{def:dv:variance-swaps-and-volatility-derivatives:gamma}{gamma swap} n'explose pas quand l'indice baisse, puisque la variance de chaque jour est multipliée par
$S_t/S_0$, qui est faible après un krach. Il est donc un produit beaucoup moins dangereux à vendre que le \omterm{def:dv:variance-swaps-and-volatility-derivatives:vs}{swap de variance}. Sur
la surface du chapitre~9, son strike à un an vaut 20.8~\%, entre le niveau à la monnaie de 19.2~\% et celui du \omterm{def:dv:variance-swaps-and-volatility-derivatives:vs}{swap de variance},
soit 22.1~\%. Un corridor descendant ($U=S_0$) isole la variance payée sur les marchés baissiers, la part du \omterm{def:dv:variance-swaps-and-volatility-derivatives:vs}{swap de variance}
que le skew rend chère.

\section{Contrats à terme et options sur l'indice de volatilité}

Les contrats à terme sur l'indice de volatilité se négocient depuis mars 2004 et les options depuis février 2006. Le chapitre~12 a montré
que l'indice est, en version idéalisée, la racine carrée de la \omterm{def:dv:rough-volatility-and-forward-variance-models:fwdvar}{variance forward} moyenne sur trente jours, et que son
contrat à terme vaut $\E[\mathrm{VIX}_T]$. Il se situe donc sous la volatilité du swap de \omterm{def:dv:rough-volatility-and-forward-variance-models:fwdvar}{variance forward}
$\sqrt{\E[\mathrm{VIX}_T^2]}$ d'un écart de convexité. L'écart est le même effet de Jensen que celui du \omterm{def:dv:variance-swaps-and-volatility-derivatives:volswap}{swap de volatilité},
appliqué à trente jours de variance vus depuis $T$.

\begin{example}[La courbe des contrats à terme sous Heston]\label{ex:dv:variance-swaps-and-volatility-derivatives:vixcurve}
Sous le \omterm{def:dv:stochastic-volatility:heston}{modèle de Heston} calibré, l'indice vaut 17.1 aujourd'hui. Le contrat à terme à trois mois vaut 18.1, contre une
volatilité forward du \omterm{def:dv:variance-swaps-and-volatility-derivatives:vs}{swap de variance} sur trente jours de 20.4, le contrat à six mois 19.3 contre 22.2, et celui à un an 20.5 contre 23.7
(\cref{fig:dv:variance-swaps-and-volatility-derivatives:vix}, à gauche). Le modèle calcule exactement les contrats à terme~:
$\mathrm{VIX}_T^2=\bar v+(v_T-\bar v)\frac{1-e^{-\kappa\Delta}}{\kappa\Delta}$ est affine en $v_T$, dont la
transformée de Laplace est connue, si bien que la même identité en racine carrée s'applique. Un desk qui lit la volatilité forward sur la
courbe des contrats à terme sans tenir compte de l'écart surestime le niveau juste des contrats à terme d'un à trois points.
\end{example}

\begin{definition}[Option sur le VIX]\label{def:dv:variance-swaps-and-volatility-derivatives:vixopt}
Une \emph{option sur le VIX}\index{option sur le VIX} est une option européenne sur l'indice de volatilité à son échéance. Elle est
réglée en espèces, et son sous-jacent naturel est le contrat à terme sur le VIX de même échéance, si bien qu'elle se cote au moyen de la
formule de Black sur ce contrat à terme.
\end{definition}

Une \omterm{def:dv:variance-swaps-and-volatility-derivatives:vixopt}{option sur le VIX} est une option sur la courbe de \omterm{def:dv:rough-volatility-and-forward-variance-models:fwdvar}{variance forward}, si bien que son smile mesure la distribution de la
volatilité future. Le smile du VIX sur le marché est orienté à la hausse~: les calls sur forte volatilité sont chers. Baldeaux et Badran
soulignent ce skew haussier observé. Sous le \omterm{def:dv:stochastic-volatility:heston}{modèle de Heston} calibré, le smile à trois mois est orienté dans l'autre sens,
allant de 110~\% à 80~\% du contrat à terme jusqu'à 94~\% à 175~\%
(\cref{fig:dv:variance-swaps-and-volatility-derivatives:vix}, à droite). Le processus en racine carrée a une queue
droite fine. Le smile du VIX sous Rough Bergomi est presque plat (chapitre~12). Ajuster la pente du smile du VIX avec le même modèle que pour
le smile de l'indice relève du problème de calibration conjointe du chapitre~12.

\begin{omfigure}
\begin{tikzpicture}
\begin{groupplot}[group style={group size=2 by 1,horizontal sep=1.6cm},
  width=6.6cm,height=5cm,
  tick label style={font=\small},label style={font=\small},grid=major,grid style={gray!25},table/col sep=comma]
\nextgroupplot[xlabel={échéance des futures (mois)},ylabel={points de volatilité},xmin=0,xmax=12,ymin=16,ymax=25,
  title={\small futures sur l'indice, Heston},
  legend columns=1,legend cell align=left,legend style={font=\footnotesize,at={(0.5,-0.32)},anchor=north,draw=none,fill=none}]
\addplot[omThm,very thick,mark=*,mark size=1.4pt] table[x=months,y=future]{figdata/derivatives/14-variance-swaps-and-volatility-derivatives/vix_curve.csv};
\addplot[omDef,thick,dashed] table[x=months,y=fwd]{figdata/derivatives/14-variance-swaps-and-volatility-derivatives/vix_curve.csv};
\legend{future,volatilité forward du swap de variance à 30 jours}
\nextgroupplot[xlabel={strike (\% du future)},ylabel={volatilité implicite (\%)},xmin=75,xmax=180,ymin=85,ymax=120,
  title={\small options sur le VIX à trois mois, Heston}]
\addplot[omThm,very thick,mark=*,mark size=1.6pt] table[x=strike,y=vol]{figdata/derivatives/14-variance-swaps-and-volatility-derivatives/vix_smile.csv};
\end{groupplot}
\end{tikzpicture}

\omcaption{L'indice de volatilité sous le \omterm{def:dv:stochastic-volatility:heston}{modèle de Heston} calibré. À gauche~: les futures se situent sous la volatilité forward
du \omterm{def:dv:variance-swaps-and-volatility-derivatives:vs}{swap de variance} d'un écart de convexité croissant avec l'échéance. À droite~: le smile à trois mois des options sur
l'indice est orienté vers le bas, à l'opposé de la pente haussière du marché. Données~: le tutoriel.}\label{fig:dv:variance-swaps-and-volatility-derivatives:vix}
\end{omfigure}

\begin{definition}[Prime de risque de variance]\label{def:dv:variance-swaps-and-volatility-derivatives:vrp}
La \emph{prime de risque de variance}\index{prime de risque de variance} est l'écart entre le strike du \omterm{def:dv:variance-swaps-and-volatility-derivatives:vs}{swap de variance}
et l'espérance de la variance réalisée sous la mesure historique, $K^2-\E^{\mathbb P}[\sigma_R^2]$ (certains auteurs
prennent le signe opposé). C'est ce que touche en moyenne le vendeur de variance pour assumer le risque de
pics de variance.
\end{definition}

Lorsque la prime est positive, un vendeur de \omterm{def:dv:variance-swaps-and-volatility-derivatives:vs}{swaps de variance} la perçoit en moyenne puis la restitue, intérêts compris,
lors des krachs, lorsque la variance réalisée dépasse plusieurs fois le strike. Sa taille, son signe et ses variations d'un sous-jacent à l'autre
sont des questions empiriques. Carr et Wu (2009) en constituent l'étude classique, avec des \omterm{def:dv:variance-swaps-and-volatility-derivatives:vs}{swaps de variance}
synthétiques, et le livre~6 les mesure. La distinction compte pour ce chapitre~: l'argument de réplication
donne le strike sous la mesure de valorisation, et ne dit rien de la rentabilité de cette vente.

\section{Skew, sauts et limites de la réplication}

La réplication tient si l'indice évolue de façon continue. Sur une journée de log-rendement $r$, la couverture (le strip
plus la position quotidienne $2/S$) rapporte $2(e^r-1-r)$ tandis que la jambe du swap verse $r^2$. L'écart est
\[
2(e^r-1-r)-r^2=\frac{r^3}3+\frac{r^4}{12}+\cdots,
\]
négligeable pour les journées ordinaires, pas pour les krachs. Une baisse de 20~\% en log coûte au vendeur couvert
$0.04-0.0375=0.0025$ de variance par unité de notionnel sur cette seule journée. Dans un modèle à sauts, le strike juste
n'est donc pas celui du strip~: $\E[\sigma_R^2]T=K_{\mathrm{strip}}^2T-\E\bigl[\sum(2(e^r-1-r)-r^2)\bigr]$. Avec des sauts
négatifs, la correction est positive, et le strip sous-évalue le \omterm{def:dv:variance-swaps-and-volatility-derivatives:vs}{swap de variance}.

\begin{example}[Le strip sous les sauts de Merton]\label{ex:dv:variance-swaps-and-volatility-derivatives:jumps}
Avec le modèle de Merton du chapitre~13 (3.16 sauts par an, d'amplitude logarithmique moyenne $-5.8\%$ et de dispersion 4.6~\%), chaque saut
laisse la couverture en déficit de $-0.000182$ de variance en moyenne ($\E[J^3]/3=-0.000188$ au premier ordre). Le strike
juste dépasse celui du strip de 0.00058 de variance par an, soit 0.18 point de volatilité. Sur 100\,000 trajectoires à trois mois,
un vendeur qui fixe son strike au niveau du strip et couvre quotidiennement perd 0.17 point en moyenne (0.18 à l'erreur de simulation près), avec un écart-type de 0.40
point, et une perte de fréquence un sur cent de 1.9 points. Dans une diffusion de même variance, la même couverture présente un
écart-type de 0.014 point (\cref{fig:dv:variance-swaps-and-volatility-derivatives:hedge}).
\end{example}

\begin{omfigure}
\begin{tikzpicture}
\begin{axis}[width=10.5cm,height=5cm,ymode=log,
  xlabel={perte du swap de variance à 3 mois vendu et couvert (points de vol.)},ylabel={prob. d'une perte supérieure},
  tick label style={font=\small},label style={font=\small},
  xmin=0,xmax=3,ymin=1e-5,ymax=1,grid=major,grid style={gray!25},
  legend columns=2,legend style={font=\footnotesize,at={(0.5,-0.3)},anchor=north,draw=none,fill=none,/tikz/every even column/.append style={column sep=8pt}},
  table/col sep=comma,unbounded coords=jump]
\addplot[omDef,very thick,dashed] table[x=loss,y=diffusion]{figdata/derivatives/14-variance-swaps-and-volatility-derivatives/hedge_tail.csv};
\addplot[omThm,very thick] table[x=loss,y=jumps]{figdata/derivatives/14-variance-swaps-and-volatility-derivatives/hedge_tail.csv};
\legend{diffusion,sauts de Merton}
\end{axis}
\end{tikzpicture}

\omcaption{Le résidu de la réplication. Un \omterm{def:dv:variance-swaps-and-volatility-derivatives:vs}{swap de variance} vendu à trois mois, dont le strike est celui du strip et couvert avec
le strip et une position quotidienne de $2/S$ sur l'indice, sur 100\,000 trajectoires~: la probabilité de perdre plus qu'un
nombre donné de points de volatilité. Dans une diffusion, aucune trajectoire ne perd plus d'un dixième de point~; avec des sauts
chaque jour de krach laisse une perte, et une trajectoire sur cent perd 1.9 points. Données~: le tutoriel.}\label{fig:dv:variance-swaps-and-volatility-derivatives:hedge}
\end{omfigure}

Trois autres limites comptent en pratique.
\begin{itemize}
\item \textbf{Les strikes viennent à manquer.} Le strip n'existe que là où des options sont cotées. Lors d'un krach, l'indice chute
sous le strike le plus bas, au-delà duquel la couverture est linéaire alors que le gain continue de croître. L'erreur de troncature
dans le \cref{tab:dv:variance-swaps-and-volatility-derivatives:errors} se transforme alors d'un biais de valorisation en perte. Martin signale que les fortes variations de cours de
2008 et 2009 ont perturbé les marchés de produits dérivés de volatilité et ont conduit le marché des \omterm{def:dv:variance-swaps-and-volatility-derivatives:vs}{swaps de variance} sur titre unique à s'assécher complètement.
\item \textbf{Caps.} Un \omterm{def:dv:variance-swaps-and-volatility-derivatives:vs}{swap de variance} peut comporter un cap sur la variance réalisée, fixé à un multiple du strike. Le
cap est un call vendu sur la variance, valorisé avec un modèle de \omterm{def:dv:stochastic-volatility:sv}{volatilité de la volatilité}, et non avec le strip.
\item \textbf{Échantillonnage discret et dividendes.} Les clôtures quotidiennes, l'absence de moyenne dans l'estimateur, et le
forward ajusté des dividendes déplacent tous le strike juste de petites quantités calculables que la
valorisation au prix de marché de la construction doit appliquer de façon cohérente.
\end{itemize}

\section{Tutoriel~: répliquer un swap de variance}\label{tut:dv:variance-swaps-and-volatility-derivatives:tut}

\begin{tutorial}
\textbf{Objectif.} Valoriser un \omterm{def:dv:variance-swaps-and-volatility-derivatives:vs}{swap de variance} à partir d'un strip sur une surface synthétique, calculer le nombre de type indice et ses
erreurs, valoriser exactement sous Heston les \omterm{def:dv:variance-swaps-and-volatility-derivatives:volswap}{swaps de volatilité} et les contrats à terme sur indice, et simuler la couverture avec et sans
sauts. \textbf{État final~:} les quatre figures, le tableau et les nombres du problème du week-end.
\begin{tutsteps}
\item \textbf{La formule de l'indice}, avec $K_0$, l'option moyennée en $K_0$ et la correction de forward~:
\omcode{code/firm/varswap/firm_varswap.py}{20}{32}{La formule de l'indice de volatilité pour une échéance.}\label{lst:dv:variance-swaps-and-volatility-derivatives:index}
\item \textbf{L'erreur de saut} de la couverture par \omterm{def:dv:variance-swaps-and-volatility-derivatives:log}{contrat log}, une ligne qui détermine le résidu de la couverture~:
\omcode{code/firm/varswap/firm_varswap.py}{74}{78}{Insuffisance de la couverture pour un rendement $r$.}\label{lst:dv:variance-swaps-and-volatility-derivatives:jump}
\item \textbf{Le \omterm{def:dv:variance-swaps-and-volatility-derivatives:volswap}{swap de volatilité} de Heston} par la transformée de Laplace de la variance intégrée et l'identité
de la racine carrée~:
\omcode{code/firm/varswap/firm_varswap.py}{92}{114}{Strike exact du swap de volatilité sous Heston.}\label{lst:dv:variance-swaps-and-volatility-derivatives:volswap}
\item \textbf{Exécuter} \texttt{dv\_varswap.index\_errors()}, \texttt{heston\_terms()}, \texttt{vix\_curve()},
\texttt{hedge\_pnl()} et \texttt{fig\_varswap.py}.
\end{tutsteps}
\textbf{Pour aller plus loin.} Remplacer l'identité de la racine carrée par l'approximation du second ordre et mesurer son
erreur à un an~; couper le strip là où le prix du put du modèle descend sous 0.05 (un «~zero bid~») et recalculer le
nombre à trente jours~; valoriser un \omterm{def:dv:variance-swaps-and-volatility-derivatives:corridor}{swap de variance corridor} à la baisse avec $U=100$.
\end{tutorial}

\section{Construction~: variance et swaps de volatilité}\label{bld:dv:variance-swaps-and-volatility-derivatives:varswap}

\begin{build}
\bfield{Rôle} Le pricer et la valorisation au prix de marché de la firme miniature pour les \omterm{def:dv:variance-swaps-and-volatility-derivatives:vs}{swaps de variance}, ainsi que les valeurs
du \omterm{def:dv:variance-swaps-and-volatility-derivatives:volswap}{swap de volatilité} et du contrat à terme sur indice auxquelles son desk de volatilité cote.
\bfield{Interface} \texttt{index\_variance(strikes, calls, puts, fwd, t, rate)}~; \texttt{strip\_from\_vols}~;
\texttt{log\_payoff\_strip}~; \texttt{realised\_variance(prices)}~; \texttt{variance\_notional(vega\_notional,
strike\_vol)}~; \texttt{mark\_to\_market(var\_notional, strike\_vol, realised\_var, t\_elapsed,
fair\_var\_remaining, t\_remaining, df)}~; \texttt{jump\_error(r)}~; \texttt{vol\_swap\_approx}~;
\texttt{heston\_mean\_variance}, \texttt{heston\_vol\_swap}, \texttt{heston\_vix\_future}.
\bfield{Règles} Les strikes sont cotés en volatilité et réglés en variance~; l'estimateur de variance réalisée est celui du
contrat (clôtures quotidiennes, 252 jours, sans moyenne)~; le strip utilise le forward et non le spot~; chaque valeur du modèle est
vérifiée par une simulation dans les tests.
\bfield{Tests d'acceptation} \texttt{code/firm/varswap/tests/}~: le nombre d'indice d'un smile plat est sa volatilité,
le forward étant entre les strikes~; le strip reproduit le payoff log à l'intérieur de son intervalle~; les conventions
(conversion du notionnel, valorisation au prix de marché)~; l'erreur de saut se comporte comme $r^3/3$~; les transformées de Laplace de Heston
valent un en zéro, le \omterm{def:dv:variance-swaps-and-volatility-derivatives:volswap}{swap de volatilité} tend vers le \omterm{def:dv:variance-swaps-and-volatility-derivatives:vs}{swap de variance} lorsque la \omterm{def:dv:stochastic-volatility:sv}{volatilité de la volatilité} s'annule et concorde avec une
simulation, et les contrats à terme sur indice se situent sous le forward.
\bfield{Pour aller plus loin} \omterm{def:dv:variance-swaps-and-volatility-derivatives:vs}{Swaps de variance} plafonnés sous Heston~; swaps corridor et swaps gamma à partir d'une chaîne cotée~; les règles
du cours spécial d'ouverture pour le règlement-livraison.
\end{build}

\begin{omsources}
\item K.~Demeterfi, E.~Derman, M.~Kamal et J.~Zou, ``More than you ever wanted to know about volatility swaps'',
Goldman Sachs Quantitative Strategies Research Notes (March 1999).
\item Cboe Global Markets, ``Cboe Volatility Index Methodology'', version 6.0 (2026).
\item P.~Carr et R.~Lee, ``Volatility derivatives'', \emph{Annual Review of Financial Economics} 1 (2009) 319--339.
\item P.~Carr et L.~Wu, ``Variance risk premiums'', \emph{Review of Financial Studies} 22(3) (2009) 1311--1341.
\item I.~Martin, ``Simple variance swaps'', NBER Working Paper 16884 (2011).
\item J.~Baldeaux et A.~Badran, ``Consistent modeling of VIX and equity derivatives using a 3/2 plus jumps
model'', arXiv 1203.5903 (2012).
\end{omsources}

\section{Exercices}

\begin{exercise}[$\star$]\label{exo:dv:variance-swaps-and-volatility-derivatives:1}
Un client veut 50\,000 de \omterm{def:dv:variance-swaps-and-volatility-derivatives:veganot}{notionnel vega} sur un \omterm{def:dv:variance-swaps-and-volatility-derivatives:vs}{swap de variance} dont le strike est fixé à 20~\%. Calculer le \omterm{def:dv:variance-swaps-and-volatility-derivatives:varnot}{notionnel de variance}, et
dire ce que le swap verse si la volatilité réalisée vaut 25~\%.
\end{exercise}

\begin{exercise}[$\star$]\label{exo:dv:variance-swaps-and-volatility-derivatives:2}
Pourquoi un \omterm{def:dv:variance-swaps-and-volatility-derivatives:vs}{swap de variance} a-t-il besoin d'une transaction quotidienne sur l'indice en plus du strip d'options~?
\end{exercise}

\begin{exercise}[$\star$]\label{exo:dv:variance-swaps-and-volatility-derivatives:3}
Montrer que le payoff répliquant du swap gamma $\frac2{S_0}(S\ln(S/S_0)-S+S_0)$ a une dérivée seconde égale à $2/(S_0S)$,
et expliquer pourquoi son strike est inférieur à celui du \omterm{def:dv:variance-swaps-and-volatility-derivatives:vs}{swap de variance} sur un indice à asymétrie baissière.
\end{exercise}

\begin{exercise}[$\star\star$]\label{exo:dv:variance-swaps-and-volatility-derivatives:4}
Après neuf mois, un \omterm{def:dv:variance-swaps-and-volatility-derivatives:vs}{swap de variance} d'un an, de strike 22~\%, a réalisé 18~\%, et la volatilité du \omterm{def:dv:variance-swaps-and-volatility-derivatives:vs}{swap de variance} à trois mois vaut 16~\%. Valoriser
le swap par unité de \omterm{def:dv:variance-swaps-and-volatility-derivatives:varnot}{notionnel de variance} (variance en pourcentage au carré, taux nuls).
\end{exercise}

\begin{exercise}[$\star\star$]\label{exo:dv:variance-swaps-and-volatility-derivatives:5}
Utiliser l'approximation du second ordre pour estimer le strike du \omterm{def:dv:variance-swaps-and-volatility-derivatives:volswap}{swap de volatilité} lorsque le strike du \omterm{def:dv:variance-swaps-and-volatility-derivatives:vs}{swap de variance} vaut 20~\%
et que l'écart type de la variance réalisée vaut 0.02.
\end{exercise}

\begin{exercise}[$\star\star$]\label{exo:dv:variance-swaps-and-volatility-derivatives:6}
Pourquoi l'ajustement de convexité de l'\cref{ex:dv:variance-swaps-and-volatility-derivatives:convexity} atteint-il son maximum vers
un an~?
\end{exercise}

\begin{exercise}[$\star\star\star$]\label{exo:dv:variance-swaps-and-volatility-derivatives:7}
\emph{Programmation.} Calculer le nombre d'indice à trente jours avec des strikes tous les 1 de 80 à 120, puis retirer les
strikes dont l'option hors de la monnaie vaut moins de 0.05. De combien le nombre baisse-t-il~?
\end{exercise}

\begin{exercise}[$\star\star\star$]\label{exo:dv:variance-swaps-and-volatility-derivatives:8}
\emph{Repérer la faille.} «~Nous vendons des \omterm{def:dv:variance-swaps-and-volatility-derivatives:vs}{swaps de variance} au prix du strip et nous détenons le strip et le delta quotidien~; le
livre est entièrement couvert, donc la prime que nous facturons au-dessus du strip est un profit pur.~»
\end{exercise}

\section{Problème~: Vol ou variance}

\begin{problem}[{Problème du week-end --- le prix d'une racine carrée}]\label{pb:dv:variance-swaps-and-volatility-derivatives:1}
Un client demande des cotations à un an sur un \omterm{def:dv:variance-swaps-and-volatility-derivatives:vs}{swap de variance} et sur un \omterm{def:dv:variance-swaps-and-volatility-derivatives:volswap}{swap de volatilité} portant sur l'indice de la surface du
chapitre~9. Le desk dispose de la surface, de la calibration de Heston du chapitre~10 et de la construction.

\medskip\noindent\textbf{Partie I --- Le \omterm{def:dv:variance-swaps-and-volatility-derivatives:vs}{swap de variance}.}
\begin{enumerate}
\item Donner le strike du \omterm{def:dv:variance-swaps-and-volatility-derivatives:vs}{swap de variance} à un an à partir du strip et de la volatilité à la monnaie.
\item Quels strikes portent le strike~? Donner la part des puts.
\item Pour un \omterm{def:dv:variance-swaps-and-volatility-derivatives:veganot}{notionnel vega} de 100\,000, donner le \omterm{def:dv:variance-swaps-and-volatility-derivatives:varnot}{notionnel de variance} et le payoff à 30~\% de volatilité réalisée.
\item De combien le nombre de type indice à trente jours perdrait-il avec un strip qui s'arrête à 95~?
\item Donner le strike du swap gamma à un an.
\end{enumerate}

\medskip\noindent\textbf{Partie II --- Le \omterm{def:dv:variance-swaps-and-volatility-derivatives:volswap}{swap de volatilité} sous Heston.}
\begin{enumerate}[resume]
\item Donner le strike du \omterm{def:dv:variance-swaps-and-volatility-derivatives:vs}{swap de variance} à un an selon Heston, et expliquer l'écart avec celui du strip.
\item Donner le strike du \omterm{def:dv:variance-swaps-and-volatility-derivatives:volswap}{swap de volatilité} à un an.
\item Donner l'ajustement de convexité à un mois, à un an et à trois ans.
\item Quel paramètre de Heston gouverne l'ajustement, et dans quel sens~?
\item Pourquoi le \omterm{def:dv:variance-swaps-and-volatility-derivatives:volswap}{swap de volatilité} ne peut-il pas être répliqué par un strip statique~?
\end{enumerate}

\medskip\noindent\textbf{Partie III --- L'indice et ses options.}
\begin{enumerate}[resume]
\item Donner le niveau de l'indice ainsi que les contrats à terme à trois mois et à un an selon Heston.
\item Donner l'écart de convexité à trois mois.
\item Décrire le smile du VIX à trois mois selon Heston.
\item Comment le desk couvrirait-il un \omterm{def:dv:variance-swaps-and-volatility-derivatives:volswap}{swap de volatilité}~?
\item Quelle prime au-dessus du strip demanderiez-vous à un client qui achète le \omterm{def:dv:variance-swaps-and-volatility-derivatives:vs}{swap de variance}, si l'indice peut sauter~?
\end{enumerate}

\medskip\noindent\textbf{Partie IV --- Jugement.}
\begin{enumerate}[resume]
\item Le client affirme que le \omterm{def:dv:variance-swaps-and-volatility-derivatives:volswap}{swap de volatilité} «~doit être évalué au strike du \omterm{def:dv:variance-swaps-and-volatility-derivatives:vs}{swap de variance}, puisque les deux sont
de la volatilité~». Répondre en deux phrases.
\item De quelle hypothèse du modèle l'ajustement de convexité dépend-il le plus~?
\item Quelle réserve faudrait-il constituer sur un \omterm{def:dv:variance-swaps-and-volatility-derivatives:volswap}{swap de volatilité} vendeur couvert par des \omterm{def:dv:variance-swaps-and-volatility-derivatives:vs}{swaps de variance}~?
\item Énoncer le \emph{résultat nommé}~: l'ajustement de convexité entre le strike à un an du \omterm{def:dv:variance-swaps-and-volatility-derivatives:vs}{swap de variance} et
celui du \omterm{def:dv:variance-swaps-and-volatility-derivatives:volswap}{swap de volatilité} sur la surface du chapitre, sous la dynamique de Heston.
\item En une phrase~: pourquoi est-ce la variance, et non la volatilité, qui est négociée~?
\end{enumerate}
\end{problem}

\section{Questions d'entretien}

\begin{interviewq}[$\star$ \iqroles{trader, chercheur}]\label{iq:dv:variance-swaps-and-volatility-derivatives:1}
Comment répliquez-vous un \omterm{def:dv:variance-swaps-and-volatility-derivatives:vs}{swap de variance}~?
\end{interviewq}

\begin{interviewq}[$\star$ \iqroles{trader}]\label{iq:dv:variance-swaps-and-volatility-derivatives:2}
Pourquoi le strike du \omterm{def:dv:variance-swaps-and-volatility-derivatives:vs}{swap de variance} se situe-t-il au-dessus de la volatilité à la monnaie d'un indice~?
\end{interviewq}

\begin{interviewq}[$\star\star$ \iqroles{chercheur}]\label{iq:dv:variance-swaps-and-volatility-derivatives:3}
Le strike juste d'un \omterm{def:dv:variance-swaps-and-volatility-derivatives:volswap}{swap de volatilité} est-il supérieur ou inférieur à celui d'un \omterm{def:dv:variance-swaps-and-volatility-derivatives:vs}{swap de variance}~? De combien, à peu près~?
\end{interviewq}

\begin{interviewq}[$\star\star$ \iqroles{développeur}]\label{iq:dv:variance-swaps-and-volatility-derivatives:4}
Implémentez la formule de l'indice de volatilité pour une échéance. Quels cas limites comptent~?
\end{interviewq}

\begin{interviewq}[$\star\star$ \iqroles{risque, trader}]\label{iq:dv:variance-swaps-and-volatility-derivatives:5}
Que se passe-t-il pour un dealer vendeur d'un \omterm{def:dv:variance-swaps-and-volatility-derivatives:vs}{swap de variance}, couvert avec le strip, lorsque l'indice chute de 20~\%~?
\end{interviewq}

\begin{interviewq}[$\star\star\star$ \iqroles{trader, risque}]\label{iq:dv:variance-swaps-and-volatility-derivatives:6}
Pourquoi un contrat à terme sur le VIX doit-il être évalué en dessous de la volatilité forward du \omterm{def:dv:variance-swaps-and-volatility-derivatives:vs}{swap de variance}, et comment négocieriez-vous l'écart~?
\end{interviewq}
